\documentclass[a4paper,12pt]{article}
\usepackage[utf8]{inputenc}
\usepackage[T1]{fontenc}
\usepackage[slovak]{babel}
\usepackage{geometry}
\usepackage{longtable}
\usepackage{array}
\usepackage{booktabs}

\geometry{margin=2.5cm}

\title{Užívateľské testovanie 1}
\date{}

\begin{document}

\maketitle

\section*{Testované časti}
Testovanie bolo zamerané na základné ovládanie aplikácie a navigáciu v nej, bez úpravy obrázkov.
\begin{itemize}
    \item Úvodná obrazovka
    \item Nastavenia
    \item Nápoveda
    \item Štatistiky
    \item Tutoriál
    \item Vloženie a export obrázka
\end{itemize}

\section*{Zameranie testovania (cieľ)}
\begin{itemize}
    \item čo ti pri tejto operácii chýbalo?
    \item čo sa stalo inak oproti očakávaniu?
    \item čo ti pri tejto operácii vadilo?
    \item čo by si zmenil?
    \item páčil sa ti dizajn?
    \item páčilo sa ti rozloženie?
\end{itemize}

\newpage
\section*{Hodnotenie použiteľnosti}
\begin{longtable}{>{\bfseries}l l}
\toprule
Hodnotenie & Popis \\
\midrule
0 & užívateľ to nenašiel vôbec \\
1 & užívateľ to našiel s veľkou nápovedou \\
2 & užívateľ to našiel s malou nápovedou \\
3 & užívateľ sa musel zamyslieť viac \\
4 & užívateľ sa musel zamyslieť \\
5 & užívateľ to našiel bez váhania \\
\bottomrule
\end{longtable}

\section*{Otázky / úlohy}

\subsection*{1. Okno nastavení}
\begin{itemize}
    \item Zmena jazyka aplikácie
    \item Zmena farebného režimu 
    \item Prepínanie ostatných nastavení 
\end{itemize}

\begin{longtable}{>{\bfseries}l c}
\toprule
Používateľ & Skóre \\
\midrule
Užívateľ 1 & 5 \\
Užívateľ 2 & 5 \\
Užívateľ 3 & 5 \\
Výsledok & 5 \\
\bottomrule
\end{longtable}

\textbf{Poznámky / zistenia:}
\begin{itemize}
    \item Divný efekt pri hover nad prepínačom jazykov (príde mi to neintuitívne)
    \item Zmena farebného režimu by mohla byť implementovaná rovnakým typom tlačidla ako ostatné veci (využiť toggle button)
    \item Nápoveda pre jazyk sa zobrazuje divne v SK a CZ
\end{itemize}

\newpage
\subsection*{2. Okno nastavení}
\begin{itemize}
    \item Aká je aktuálna verzia aplikácie 
    \item Zisti dátum vydania konkrétnej verzie aplikácie 
\end{itemize}

\begin{longtable}{>{\bfseries}l c}
\toprule
Používateľ & Skóre \\
\midrule
Užívateľ 1 & 4 \\
Užívateľ 2 & 5 \\
Užívateľ 3 & 5 \\
Výsledok & 4.7 \\
\bottomrule
\end{longtable}

\textbf{Poznámky / zistenia:}
\begin{itemize}
    \item Nezaujíma ma to, ale nie je to výrazné, takže to asi neprekáža
\end{itemize}

\subsection*{3. Nápoveda, tutoriál}
\begin{itemize}
    \item Spusti tutoriál
\end{itemize}

\begin{longtable}{>{\bfseries}l c}
\toprule
Používateľ & Skóre \\
\midrule
Užívateľ 1 & 4 \\
Užívateľ 2 & 4 \\
Užívateľ 3 & 5 \\
Výsledok & 4.3 \\
\bottomrule
\end{longtable}

\textbf{Poznámky / zistenia:}
\begin{itemize}
    \item Kratší popis, len tak na 2 riadky a mať možnosť to rozkliknúť pre podrobnejší popis
    \item Možnosť otvoriť pravý panel (panel s nastavením nástrojov) počas tutoriálu
    \item Dĺžka popisu je vyhovujúca
    \item Zvýraznenie len tých častí, ktorých sa to týka
    \item Zvýraznenie by mohlo mať zaoblené rohy, aby to ladilo so zvyškom
    \item Export je zahrnutý v 2 krokoch tutoriálu
\end{itemize}

\newpage
\subsection*{4. Nápoveda, štatistiky}
\begin{itemize}
    \item Zobraz štatistiky
\end{itemize}

\begin{longtable}{>{\bfseries}l c}
\toprule
Používateľ & Skóre \\
\midrule
Užívateľ 1 & 5 \\
Užívateľ 2 & 4 \\
Užívateľ 3 & 4 \\
Výsledok & 4.3 \\
\bottomrule
\end{longtable}

\textbf{Poznámky / zistenia:}
\begin{itemize}
    \item Hapruje zobrazenie správneho jazyku pri prepínaní (to isté pri dark/light režime)
\end{itemize}

\subsection*{5. Privacy and data}
\begin{itemize}
    \item Vymaž užívateľské preferencie aplikácie 
\end{itemize}

\begin{longtable}{>{\bfseries}l c}
\toprule
Používateľ & Skóre \\
\midrule
Užívateľ 1 & 3 \\
Užívateľ 2 & 3 \\
Užívateľ 3 & 2 \\
Výsledok & 2.7 \\
\bottomrule
\end{longtable}

\textbf{Poznámky / zistenia:}
\begin{itemize}
    \item Skôr to spraviť ako tlačidlo, takto to vyzerá ako odkaz
\end{itemize}

\subsection*{6. Nápoveda – skratky a obmedzenia}
\begin{itemize}
    \item Nájdi vysvetlenie skratiek
    \item Nájdi zoznam nepodporovaných funkcií/obmedzení
\end{itemize}

\begin{longtable}{>{\bfseries}l c}
\toprule
Používateľ & Skóre \\
\midrule
Užívateľ 1 & 5 \\
Užívateľ 2 & 5 \\
Užívateľ 3 & 5 \\
Výsledok & 5 \\
\bottomrule
\end{longtable}

\textbf{Poznámky / zistenia:}
\begin{itemize}
    \item /-
\end{itemize}

\subsection*{7. Upload a export obrázku}
\begin{itemize}
    \item Otvor obrázok z počítača
    \item Otvor obrázok iným spôsobom ako oknom pre prehľadávanie súborov
    \item Exportuj obrázok 
\end{itemize}

\begin{longtable}{>{\bfseries}l c}
\toprule
Používateľ & Skóre \\
\midrule
Užívateľ 1 & 5 \\
Užívateľ 2 & 5 \\
Užívateľ 3 & 5 \\
Výsledok & 5 \\
\bottomrule
\end{longtable}

\textbf{Poznámky / zistenia:}
\begin{itemize}
    \item /-
\end{itemize}

\subsection*{8. Zatvorenie súborov}
\begin{itemize}
    \item Otvor viac ako jeden obrázok a zatvor všetky súčasne 
\end{itemize}

\begin{longtable}{>{\bfseries}l c}
\toprule
Používateľ & Skóre \\
\midrule
Užívateľ 1 & 4 \\
Užívateľ 2 & 4 \\
Užívateľ 3 & 4 \\
Výsledok & 4 \\
\bottomrule
\end{longtable}

\textbf{Poznámky / zistenia:}
\begin{itemize}
    \item Všetci to hľadali, ale nakoniec to našli pri pokuse zavrieť len jeden obrázok
\end{itemize}

\section*{Globálne hodnotenie}

\begin{longtable}{l c}
\toprule
Úloha / Otázka & Priemer / Výsledok \\
\midrule
1. Okno nastavení – jazyk a farebný režim & 5.0 \\
2. Okno nastavení – verzia aplikácie & 4.7 \\
3. Nápoveda, tutoriál & 4.3 \\
4. Nápoveda, štatistiky & 4.3 \\
5. Privacy and data & 2.7 \\
6. Nápoveda – skratky a obmedzenia & 5.0 \\
7. Upload a export obrázku & 5.0 \\
8. Zatvorenie súborov & 4.0 \\
\bottomrule
\end{longtable}

\textbf{Celkové priemerné hodnotenie všetkých úloh:} 4.4 / 5

\end{document}
